\chapter*{Introduction}\label{chap:intro}

This text presents an introduction to algebra suitable for upper-level
undergraduate or beginning graduate courses. While there is a very
extensive offering of textbooks at this level, in my experience
teaching this material I have invariably felt the need for a
self-contained text that would start `from zero' (in the sense of not
assuming that the reader has had substantial previous exposure to the
subject) but that would impart from the very beginning a rather
modern, categorically minded viewpoint and aim at reaching a good
level of depth. Many textbooks in algebra brilliantly satisfy some,
but not all, of these requirements. This book is my attempt at
providing a working alternative.  There is a widespread perception
that categories should be avoided at first blush, that the abstract
language of categories should not be introduced until a student has
toiled for a few semesters through example-driven illustrations of the
nature of a subject like algebra. According to this viewpoint,
categories are only tangentially relevant to the main topics covered
in a beginning course, so they can simply be mentioned occasionally
for the general edification of the reader, who will in time learn
about them (by osmosis?). Paraphrasing a reviewer of a draft of the
present text, `Discussions of categories at this level are the reason
why God created appendices'.  It will be clear from a cursory glance
at the table of contents that I think otherwise. In this text,
categories are introduced on page 18, after a scant reminder of the
basic language of naive set theory, for the main purpose of providing
a context for universal properties. These are in turn evoked
constantly as basic definitions are introduced. The word `universal'
appears at least 100 times in the first three chapters.  I believe
that awareness of the categorical language, and especially some
appreciation of universal properties, is particularly helpful in
approaching a subject such as algebra `from the beginning'. The reader
I have in mind is someone who has reached a certain level of
mathematical maturity---for example, who needs no special assistance
in grasping an induction argument---but may have only been exposed to
algebra in a very cursory manner. My experience is that many
upper-level undergraduates or beginning graduate students at Florida
State University and at comparable institutions fit this
description. For these students, seeing the many introductory concepts
in algebra as instances of a few powerful ideas (encapsulated in
suitable universal properties) helps to build a comforting unifying
context for these notions. The amount of categorical language needed
for this catalyzing function is very limited; for example, functors
are not really necessary in this acclimatizing stage.  Thus, in my
mind the benefit of this approach is precisely that it helps a true
beginner, if it is applied with due care. This is my experience in the
classroom, and it is the main characteristic feature of this text. The
very little categorical language introduced at the outset informs the
first part of the book, introducing in general terms groups, rings,
and modules. This is followed by a (rather traditional) treatment of
standard topics such as Sylow theorems, unique factorization,
elementary linear algebra, and field theory. The last third of the
book wades into somewhat deeper waters, dealing with tensor products
and Hom (including a first introduction to Tor and Ext) and including
a final chapter devoted to homological algebra. Some familiarity with
categorical language appears indispensable to me in order to
appreciate this latter material, and this is hopefully
uncontroversial.  Having developed a feel for this language in the
earlier parts of the book, students find the transition into these
more advanced topics particularly smooth.  A first version of this
book was essentially a careful transcript of my lectures in a run of
the (three-semester) algebra sequence at FSU.\@ The chapter on
homological algebra was added at the instigation of Ed Dunne, as were
a very substantial number of the exercises. The main body of the text
has remained very close to the original `transcript' version: I have
resisted the temptation of expanding the material when revising it for
publication. I believe that an effective introductory textbook (this
is Chapter 0, after all. . . ) should be realistic: it must be
possible to cover in class what is covered in the book. Otherwise, the
book veers into the `reference' category; I never meant to write a
reference book in algebra, and it would be futile (of me) to try to
ameliorate excellent available references such as Lang's `Algebra'~%
\cite{Lang2002Algebra}.  The problem sets will give an opportunity to
a teacher, or any motivated reader, to get quite a bit beyond what is
covered in the main text. To guide in the choice of exercises, I have
marked with a \exerused{} those problems that are directly referenced
from the text, and with a \lnot{} those problems that are referenced
from other problems. A minimalist teacher may simply assign all and
only the \exerused{} problems; these do nothing more than anchor the
understanding by practice and may be all that a student can
realistically be expected to work out while juggling TA duties and two
or three other courses of similar intensity as this one. The main body
of the text, together with these exercises, forms a self-contained
presentation of essential material.  The other exercises, and
especially the threads traced by those marked with \lnot{} , will
offer the opportunity to cover other topics, which some may well
consider just as essen- tial: the modular group, quaternions,
nilpotent groups, Artinian rings, the Jacobson radical, localization,
Lagrange's theorem on four squares, projective space and
Grassmannians, Nakayama's lemma, associated primes, the spectral
theorem for normal operators, etc., are some examples of topics that
make their appearance in the exercises. Often a topic is presented
over the course of several exercises, placed in appropriate sections
of the book. For example, `Wedderburn's little theorem' is mentioned
in \Cref{rem:rings:wedderburn} (that is: Remark 1.16 in Chapter III);
particular cases are presented in
\Cref{exer:rings:division-ring-p2-commutative,exer:groups2:division-ring-ord-64-comm},
and the reader eventually obtains a proof in
\Cref{exer:fields:wedderburn-thm-witts-proof}, following preliminaries
given in
\Cref{exer:fields:phi-n-a-divides-a-minus-1-implies-n-1,exer:fields:phi-n-a-divides-quotient}. The
\lnot{} label and perusal of the index should facilitate the
navigation of such topics. To help further in this process, I have
decorated every exercise with a list (added in square brackets) of the
places in the book that refer to it.  For example, an instructor
evaluating whether to assign
\Cref{exer:irred:euclid-variation-poly-values} will be immediately
aware that this exercise is quoted in
\Cref{exer:fields:infinitely-many-primes-cong-1-mod-n}, proving a
particular case of Dirichlet's theorem on primes in arithmetic
progressions, and that this will in turn be quoted in
\S{}\ref{subsec:fields:abelian-groups-as-galois-groups-over-q},
discussing the realization of abelian groups as Galois groups over Q.
I have put a high priority on the requirement that this should be a
self-contained text which essentially crosses all t's and dots all
i's, and does not require that the reader have access to other texts
while working through it. I have therefore made a conscious effort to
not quote other references: I have avoided as much as possible the
exquisitely tempting escape route `For a proof, see . . . .' This is
the main reason why this book is as thick as it is, even if so many
topics are not covered in it. Among these, commutative algebra and
representation theory are perhaps the most glaring omissions. The
first is represented to the extent of the standard basic definitions,
which allow me to sprinkle a little algebraic geometry here and there
(for example, see
\S{}\ref{sec:fields:algebraic-closure-nullstellensatz-and-a-little-algebraic}),
and of a few slightly more advanced topics in the exercises, but I
stopped short of covering, e.g., primary decompositions. The second is
missing altogether.  It is my hope to complement this book with a
`Chapter 1' in an undetermined future, where I will make amends for
these and other shortcomings.  By its nature, this book should be
quite suitable for self-study: readers working on their own will find
here a self-contained starting point which should work well as a
prelude to future, more intensive, explorations. Such readers may be
helped by the following `9-fold way' diagram of logical
interdependence of the chapters: This may however better reflect my
original intention than the final product.  For a more objective
gauge, this alternative diagram captures the web of references from a
chapter to earlier chapters, with the thickness of the lines
representing (roughly) the number of references: With the
self-studying reader especially in mind, I have put extra effort into
providing an extensive index. It is not realistic to make a fanfare
for each and every new term introduced in a text of this size by an
official `definition'; the index should help a lone traveler find the
way back to the source of unfamiliar terminology.  Internal references
are handled in a hopefully transparent way. For example,
\Cref{rem:rings:wedderburn} refers to Remark 1.16 in Chapter III; if
the reference is made from within Chapter III, the same item is called
Remark 1.16. The list in brackets following an exercise indicates
other exercises or sections in the book referring to that
exercise. For example, \Cref{exer:prelim:op-categ-def} is followed by
[\ref{exer:prelim-init-of-final-op}, \S{}\ref{subsec:linrep:functors},
\S{}\ref{subsec:homol:additive-categories},
\ref{exer:homol:opposite-of-abelian-is-abelian}]: this alerts the
reader that there are references to this problem in
\Cref{exer:prelim-init-of-final-op}, section 1.1 in Chapter VIII,
section 1.2 in Chapter IX, and
\Cref{exer:homol:opposite-of-abelian-is-abelian} (and nowhere else).
Acknowledgments. My
debt to Lang's book, to David Dummit and Richard Foote's `Abstract
Algebra,' or to Artin's `Algebra' will be evident to anyone who is
familiar with these sources. The chapter on homological algebra owes
much to David Eisenbud's appendix on the topic in his `Commutative
Algebra', to Gelfand- Manin's `Methods of homological algebra', and to
Weibel's `An introduction to homological algebra'. But in most cases
it would simply be impossible for me to retrace the original source of
an expository idea, of a proof, of an exercise, or of a specific
pedagogical emphasis: these are all likely offsprings of ideas from
any one of these and other influential references and often of
associations triggered by following the manifold strands of the World
Wide Web. This is another reason why, in a spirit of equanimity, I
resolved to essentially avoid references altogether.  In any case, I
believe all the material I have presented here is standard, and I only
retain absolute ownership of every error left in the end product.  I
am very grateful to my students for the constant feedback that led me
to write this book in this particular way and who contributed
essentially to its success in my classes. Some of the students
provided me with extensive lists of typos and outright mistakes, and I
would especially like to thank Kevin Meek, Jay Stryker, and Yong Jae
Cha for their particularly helpful comments. I had the opportunity to
try out the material on homological algebra in a course given at
Caltech in the fall of 2008, while on a sabbatical from FSU, and I
would like to thank Caltech and the audience of the course for their
hospitality and the friendly atmosphere.  Thanks are also due to MSRI
for hospitality during the winter of 2009, when the last fine-tuning
of the text was performed.  A few people spotted big and small
mistakes in preliminary versions of this book, and I will mention
Georges Elencwajg, Xia Liao, and Mirroslav Yotov for particularly
precious contributions. I also commend Arlene O'Sean and the staff at
the AMS for the excellent copyediting and production work.  Special
thanks go to Ettore Aldrovandi for expert advice, to Matilde Marcolli
for her encouragement and indispensable help, and to Ed Dunne for
suggestions that had a great impact in shaping the final version of
this book.  Support from the Max-Planck-Institut in Bonn, from the
NSA, and from Caltech, at different stages of the preparation of this
book, is gratefully acknowledged.

%%% Local Variables:
%%% mode: LaTeX
%%% TeX-engine: luatex
%%% TeX-master: "../master"
%%% End:
